%%
%% This is file `sample-sigplan.tex',
%% generated with the docstrip utility.
%%
%% The original source files were:
%%
%% samples.dtx  (with options: `all,proceedings,bibtex,sigplan')
%% 
%% IMPORTANT NOTICE:
%% 
%% For the copyright see the source file.
%% 
%% Any modified versions of this file must be renamed
%% with new filenames distinct from sample-sigplan.tex.
%% 
%% For distribution of the original source see the terms
%% for copying and modification in the file samples.dtx.
%% 
%% This generated file may be distributed as long as the
%% original source files, as listed above, are part of the
%% same distribution. (The sources need not necessarily be
%% in the same archive or directory.)
%%
%%
%% Commands for TeXCount
%TC:macro \cite [option:text,text]
%TC:macro \citep [option:text,text]
%TC:macro \citet [option:text,text]
%TC:envir table 0 1
%TC:envir table* 0 1
%TC:envir tabular [ignore] word
%TC:envir displaymath 0 word
%TC:envir math 0 word
%TC:envir comment 0 0
%%
%%
%% The first command in your LaTeX source must be the \documentclass
%% command.
%%
%% For submission and review of your manuscript please change the
%% command to \documentclass[manuscript, screen, review]{acmart}.
%%
%% When submitting camera ready or to TAPS, please change the command
%% to \documentclass[sigconf]{acmart} or whichever template is required
%% for your publication.
%%
%%
\documentclass[sigconf, anonymous]{acmart}


%%
%% \BibTeX command to typeset BibTeX logo in the docs
\AtBeginDocument{%
  \providecommand\BibTeX{{%
    Bib\TeX}}}

%% Rights management information.  This information is sent to you
%% when you complete the rights form.  These commands have SAMPLE
%% values in them; it is your responsibility as an author to replace
%% the commands and values with those provided to you when you
%% complete the rights form.
\setcopyright{acmlicensed}
\copyrightyear{2018}
\acmYear{2018}
\acmDOI{XXXXXXX.XXXXXXX}

%% These commands are for a PROCEEDINGS abstract or paper.
\acmConference[Conference acronym 'XX]{Make sure to enter the correct
  conference title from your rights confirmation emai}{June 03--05,
  2018}{Woodstock, NY}
%%
%%  Uncomment \acmBooktitle if the title of the proceedings is different
%%  from ``Proceedings of ...''!
%%
%%\acmBooktitle{Woodstock '18: ACM Symposium on Neural Gaze Detection,
%%  June 03--05, 2018, Woodstock, NY}
\acmISBN{978-1-4503-XXXX-X/18/06}

%% get rid of this for the final version
\usepackage{todonotes}
%\usepackage{listingsutf8}
\usepackage{listings}
\setlength {\marginparwidth }{2cm} %This is for todod
\usepackage{tikz}
\usetikzlibrary{positioning}
\usepackage{cryptocode}
\usepackage{lstcoq}
\usepackage[utf8]{inputenc}
\usepackage{amsmath}
\usepackage{seqsplit}
\usepackage{hyperref}
%% End of get rid of this for the final version


%%
%% Submission ID.
%% Use this when submitting an article to a sponsored event. You'll
%% receive a unique submission ID from the organizers
%% of the event, and this ID should be used as the parameter to this command.
%%\acmSubmissionID{123-A56-BU3}

%%
%% For managing citations, it is recommended to use bibliography
%% files in BibTeX format.
%%
%% You can then either use BibTeX with the ACM-Reference-Format style,
%% or BibLaTeX with the acmnumeric or acmauthoryear sytles, that include
%% support for advanced citation of software artefact from the
%% biblatex-software package, also separately available on CTAN.
%%
%% Look at the sample-*-biblatex.tex files for templates showcasing
%% the biblatex styles.
%%

%%
%% The majority of ACM publications use numbered citations and
%% references.  The command \citestyle{authoryear} switches to the
%% "author year" style.
%%
%% If you are preparing content for an event
%% sponsored by ACM SIGGRAPH, you must use the "author year" style of
%% citations and references.
%% Uncommenting
%% the next command will enable that style.
%%\citestyle{acmauthoryear}


%%
%% end of the preamble, start of the body of the document source.
\begin{document}

%%
%% The "title" command has an optional parameter,
%% allowing the author to define a "short title" to be used in page headers.
\title{One Proof, Many Implementations of Certified $\Sigma$-Protocols}

%%
%% The "author" command and its associated commands are used to define
%% the authors and their affiliations.
%% Of note is the shared affiliation of the first two authors, and the
%% "authornote" and "authornotemark" commands
%% used to denote shared contribution to the research.


%%
%% By default, the full list of authors will be used in the page
%% headers. Often, this list is too long, and will overlap
%% other information printed in the page headers. This command allows
%% the author to define a more concise list
%% of authors' names for this purpose.

%%
%% The abstract is a short summary of the work to be presented in the
%% article.
\begin{abstract}
  $\Sigma$-protocols are a core building block in a wide range of privacy-preserving applications, but 
  one particular application where they have found a widespread adoption is electronic voting.
  They remain a popular choice in electronic voting to maintain privacy 
  and verifiability of elections and are still used in prominent electronic voting systems such as 
  Helios, Belenios, and SwissPost. However, existing implementations of $\Sigma$-protocols have two 
  critical shortcomings: (i) the same code is written twice, once for the frontend (JavaScript/TypeScript)
  and once for the backend (Python/OCaml/Java), which is error-prone and may lead to inconsistencies and bugs, and (ii) none of these implementations provide formal security guarantees.
  
  In this paper, we present a certified formalisation of a comprehensive library of $\Sigma$-protocols 
  in the Rocq theorem prover, and we formally prove (special) soundness, completeness, and 
  (special honest-verifier) zero-knowledge properties. Our generic formalisation over abstract vector spaces enables us to uniformly express and compose a wide class of discrete-logarithm-based $\Sigma$-protocols, while remaining independent of any concrete group or cryptographic instantiation. Moreover, the constructive nature of our formalisation allow us to extract and compile to executable implementations in OCaml and WebAssembly and thereby
  eliminating code duplication. We demonstrate the practical applicability of our framework by implementing a voting client and server for Approval Voting, as well as an International Association for Cryptologic Research (IACR) election verifier. Our certified library provides a foundation for 
  building formally verified privacy-preserving applications. 


 


\end{abstract}


\begin{CCSXML}
<ccs2012>
  <concept>
      <concept_id>10002978.10002986.10002990</concept_id>
      <concept_desc>Security and privacy~Logic and verification</concept_desc>
      <concept_significance>500</concept_significance>
      </concept>
</ccs2012>
\end{CCSXML}
  
\ccsdesc[500]{Security and privacy~Logic and verification}

%%
%% Keywords. The author(s) should pick words that accurately describe
%% the work being presented. Separate the keywords with commas.
\keywords{Sigma Protocol, Formal Verification, Rocq, Cryptography, Zero-Knowledge Proof}



%%
%% This command processes the author and affiliation and title
%% information and builds the first part of the formatted document.
\maketitle




\section{Introduction}
Electronic voting offers convenience and efficiency in electoral processes, potentially improving 
accessibility for disabled voters, increasing participation among overseas voters, and lowering 
the costs associated with conducting elections. Moreover,
they can streamline the voting process, reduce administrative burdens, and enable 
faster tallying of results. However, electronic voting also brings significant 
challenges in ensuring the privacy and verifiability, a must for free-and-fair elections, of the voting process. Therefore, 
privacy and verifiability are two crucial elements that a voting system must ensure.
In a paper-ballot election, privacy is maintained by 
ensuring that no identifying information is present on a ballot that links it to a particular voter, 
and verifiability is achieved through scrutineers closely observing various electoral processes, 
including the counting of ballots. In order to have a similar effect in electronic voting, encryption 
is employed to safeguard the privacy of voter choices, while publicly verifiable evidence 
(such as zero-knowledge proofs) allows voters and independent observers to verify the integrity of 
the voting process. However, no matter how advanced the cryptographic techniques, they are useless 
if their software implementation contains bugs. In fact, 
the Swiss Federal Chancellery suspended electronic voting in 2019 because of 
a fatal software bug in the implementation of a cryptographic primitive 
that allowed someone to change the cast votes without being detected \cite{9152765}.
It only resumed in 2023 after getting feedbacks from academic 
community \cite{swiss_evoting_chronik} about the various components used in 
the electronic voting software. However, it is not an isolated  
instance where bugs have been found in a voting software, e.g., 
Voatz \cite{255334} (used in West Virginia, USA election),
Democracy Live Online Voting System \cite{263858} 
(used in Delaware, West Virginia, and New Jersey, USA election), 
Moscow Internet Voting System \cite{10.1007/978-3-030-51280-4_3}
(used in Russia election), and iVote System \cite{10.1007/978-3-319-22270-7_3, 10.1145/3014812.3014837} 
(used in New South Wales, Australia election). 
These unfortunate situations can largely be attributed to two factors: 
(i) closed source code (proprietary artifacts) because most of these bugs were 
discovered by researchers who inspected the code after they were made public, and 
(ii) software testing which is inadequate to rule out all the bugs 
from a software. Therefore, to avoid the situations like this, 
an electronic voting system must adhere to \textit{software independence} \cite{rivest2008notion}:
\begin{quote}
  A voting system is software-independent if an
  (undetected) change or error in its software cannot
  cause an undetectable change or error in an
  election outcome. 
\end{quote}
\noindent
Therefore, many electronic voting systems follow the principles 
of \textit{end-to-end verifiability}. End-to-end verifiability is 
defined as: 
\begin{itemize}
  \item \textit{cast-as-intended}: voters get proof that the electronic ballot accurately reflects their choices.
  \item \textit{collected-as-cast}: the election authority provides proof that the electronic ballots are received without tampering.
  \item \textit{counted-as-collected}: the election authority provides proof that the results are based on correctly counting all received ballots.
\end{itemize} 


Finally, voters and scrutineers inspect these proofs and accept or reject the election results based on 
the verification of these proofs. In practice, this means that during the execution of an election, 
electronic voting software generates data, and voters and scrutineers verify the conformity of 
this data using independently written computer programs. The rationale is that if there is 
any discrepancy in the electronic voting system, at least one scrutineer or voter would detect it.
It is important to note that end-to-end verifiability is a valuable feature because it 
allows for the verification of the election results by checking the proofs of correctness 
once the election has concluded. This ensures that the election process was conducted 
properly and that the results are accurate, but it is only in action during or after the voting has taken place.
However, can we use it as a guiding principle to develop an electronic voting software way before 
an election? The answer is yes; we can use a formal language to express end-to-end verifiability, or parts of 
it, and develop various components of an electronic voting 
software in this formal language. This process is called formal verification, or correct-by-construction,
and it ensures that the software always produces valid proofs that 
are accepted by voters and scrutineers every time. In other words, formal verification
can be employed to develop reliable electronic voting software.
In fact, formal verification has emerged as a promising solution 
to eliminate potential bugs from software, and in recent years it has been increasingly 
adopted in real-world software applications \cite{10.1145/1111037.1111042}, 
including cryptography \cite{8835346,10.1145/3133956.3134043,190894,10.1145/2701415,
10.1145/2660267.2660370,10.1145/3319535.3363211}. Given the critical reliance of end-to-end verifiability on cryptographic proofs, the software components responsible for generating and checking these proofs constitute a particularly high-value target for formal verification. Among these components, $\Sigma$-protocols stand out as the workhorse of privacy-preserving election infrastructure. 
They are used not only to prove that individual ballots are well-formed but also to underpin the zero-knowledge assurances required in mix-nets and distributed decryption. 
Consequently, a certified, bug-free implementation of $\Sigma$-protocols directly addresses the software-independence requirement by ensuring that the verifiable evidence presented to scrutineers is generated correctly by construction.  Motivated by this need, we present a formally verified library for $\Sigma$-protocols and their compositions within the Rocq proof assistant.

$\Sigma$-protocols are a fundamental class of zero-knowledge proof systems that enable a prover to convince a verifier of the correctness of a statement without revealing any additional information beyond its validity. Owing to their simplicity and security guarantees, $\Sigma$-protocols have become a core building block in a wide range of privacy-preserving applications such as anonymous credentials \cite{10.1145/2660267.2660328}, 
password-authenticated key exchange \cite{10.5555/2022815.2022838}, 
signature \cite{10.1007/0-387-34805-0_22}, threshold signature \cite{10.1007/978-3-030-81652-0_2}, etc. In the context of electronic voting, $\Sigma$-protocols play a particularly central role: they are used to prove that ballots are well-formed, that encryptions encode valid votes, that mix-nets correctly shuffle encrypted ballots, and that tallies are computed honestly, all without compromising voter privacy. These proofs underpin end-to-end verifiability by allowing voters and independent observers to publicly verify the integrity of each stage of the election while preserving ballot secrecy. As a result, the correctness and trustworthiness of $\Sigma$-protocol implementations directly impact the security of modern electronic voting systems which makes them a natural target for rigorous formal verification. In this paper, we present a certified implementation of $\Sigma$-protocols and their compositions in 
the formal language of the Rocq proof assistant. Our formalisation encompasses the basic Schnorr protocol, 
Chaum-Pedersen, Okamoto, Pedersen linear relation, encryption proofs, and compositional constructions including Parallel-$\Sigma$, 
And-$\Sigma$, Or-$\Sigma$, Eq-$\Sigma$, and Neq-$\Sigma$ protocols. For each protocol, we formally prove 
completeness, special soundness, and special honest-verifier zero-knowledge properties.
Our formalisation is developed 
generically over abstract vector spaces, enabling uniform expression and composition of discrete-logarithm-based 
$\Sigma$-protocols while remaining independent of any concrete group or cryptographic instantiation. Additionally, we have implemented efficient arithmetic procedures for prime fields and Schnorr groups (including the Helios parameters, 
with formally verified primality proofs for 2048-bit prime $p$ and 256-bit prime $q$) to 
ensure our library is practically usable. Finally, we extract OCaml and WebAssembly code from our certified 
Rocq implementation. We demonstrate the applicability of our framework by implementing a voting client, extracted to WebAssembly from Rocq proofs, and a server, extracted to OCaml from Rocq proofs, for Approval Voting and a complete Helios election verifier that validates IACR elections (2023, 2024), verifying ballot proofs, tallying, and trustee decryption proofs. To the best of our knowledge, this is the first formal verification of $\Sigma$-protocols in Rocq that supports simultaneous extraction to both WebAssembly and OCaml. Our proofs can be accessed from this \href{https://anonymous.4open.science/r/SigmaProtocol-EBA5/README.md}{\textcolor{blue}{anonymous repo}}, which contains 30,000 lines of Rocq proofs.

\subsection{Outline}
The remainder of the paper is organised as follows. 
Section \ref{sigma_protocol} introduces zero-knowledge proofs and $\Sigma$-protocols, with emphasis on their role in electronic voting. 
Section \ref{coq_theorem} briefly presents the Rocq theorem prover and the infrastructure used for extraction to targets such as WebAssembly. 
Section \ref{schn_proof} develops our formal encoding of the Schnorr protocol, our model of uniform distributions, the main $\Sigma$-protocol compositions, concrete modular arithmetic instantiations, and our approach to extractor/simulator complexity. 
Section \ref{case_studies} demonstrates the practical use of the library through an Approval Voting client-server workflow and a certified Helios election verifier for IACR data. Finally, Section \ref{rel_work} discusses related work, current limitations, and directions for future research.



\section{$\Sigma$-Protocols}\label{sigma_protocol}
Despite the recent growth and popularity of zero-knowledge succinct 
arguments of knowledge (zk-SNARKs), $\Sigma$-protocols remain a
leading cryptographic proof system in privacy-preserving voting systems. 
$\Sigma$-protocols are used in Helios (a widely popular voting system) \cite{adida2008helios}, 
Belenios (used in French elections for overseas 
voters\footnote{VVFE, used in French elections, is a derivative of (part of) Belenios and 
inherited all the zero-knowledge proofs from Belenios. \url{https://gitlab.inria.fr/vvfe/vvfe}}), \cite{cortier2023french}, and 
SwissPost (used in Swiss elections for overseas voters) voting 
systems \cite{10.1007/978-3-031-15911-4_4}. 
Moreover, recently the zero-knowledge proof community started a movement to standardise
$\Sigma$-protocols \cite{ZKProof}.
One of the main reasons for its popularity in privacy-preserving voting systems is its
efficiency and well-understood security.


%\textbf{Introduce here that sigma protocol is efficient and simple and all other details about sigma protocol.} 
Zero-knowledge proofs were first studied by Goldwasser, Micali, and Rackoff \cite{10.1145/22145.22178} and 
are possible for all problems in $NP$ \cite{10.1145/116825.116852}. Recall that $NP$ is the class of 
problems that have efficient verifiers; that is, for a given problem in the $NP$ class, there is 
a polynomial-time algorithm that can verify whether a given solution, or witness, to the problem is correct or not.
More formally, for a given binary $NP$-relation $R$ and a statement $x$, zero-knowledge proof allows a prover 
$P$ to convince a verifier $V$ that they posses a witness $w$--polynomial in the length of $x$--such 
that ($x$, $w$) $\in$ $R$. To achieve this, $P$ and $V$ engage in an 
interactive protocol and at the end $V$ either accepts or rejects the proof.
A $\Sigma$-protocol is an efficient 3-move zero-knowledge proof (where the 
verifier is assumed to be honest) characterised by completeness, (special) soundness, and 
(special honest-verifier) zero-knowledge. More formally: (i) \textbf{completeness:}
  when $P$ and $V$ follow the protocol, $V$ always accepts. 
  In other words, if $P$ knows $w$ for $x$ such that ($x$, $w$) $\in$ $R$,  
  then $P$ can successfully construct a proof $(a, c, r)$ that always passes the verification equation, 
  i.e., $V$ accepts; (ii) \textbf{special soundness:}
there exists a probabilistic polynomial-time extractor $E$ which given 
two accepting transcripts $(a, c, r)$ and $(a, c', r')$ with $c\neq c'$ for $x$, 
it can extract a witness $w$ such that $(x, w) \in R$; and (iii) \textbf{special honest-verifier zero-knowledge:}
there exists a probabilistic polynomial-time simulator \(S\) such that, for every \(x \in L_R\) and every challenge \(c \in C\), \(S\) outputs a transcript \((a,c,r)\) whose distribution is identical to that of a real interaction between an honest prover \(P\) (holding a witness \(w\) with \((x,w)\in R\)) and an honest verifier \(V\) using challenge $c$. Furthermore, for any \(x \notin L_R\), the simulator \(S\) is only required to output an arbitrary accepting transcript \((a,c,r)\) for the given challenge \(c \in C\).


The notion of a $\Sigma$-protocol was first identified and studied by Cramer~\cite{cramer1996modular},
as an appropriate abstraction of the (special) honest-verifier zero-knowledge, special-sound protocols considered in~\cite{cramer1994proofs}.
Probably the best-known $\Sigma$-protocol is due to Schnorr~\cite{schnorr1991efficient}.
The Schnorr protocol is used extensively in privacy-preserving voting systems, e.g, Helios, Belenios, SwissPost, etc. Therefore, 
our formalisation is geared towards this protocol. Moreover, we formalise the $n$-ary composition of Schnorr $\Sigma$-protocols, instantiating Parallel, And, Or, Eq, and Neq constructions for arbitrary $n$. In addition, we also formalise the Okamoto protocol, encryption proofs (demonstrating that ElGamal ciphertexts encrypt specific values), decryption proofs (verifying correct decryption by election trustees), and protocol for proving knowledge of committed values satisfying linear constraints~\cite{Bra97}. Additionally, we instantiate Eq-composition to obtain the Chaum-Pedersen protocol for proving equality of discrete logarithms. Ideally, given the abstraction of our formalisation (\textit{vector space}), we could have 
abstracted $\Sigma$-protocols and their composition in a generic way \cite{10.1007/978-3-642-02384-2_17},
but we decided to focus on the Schnorr protocol and its composition to ensure familiar APIs 
for electronic voting in WebAssembly and OCaml.




\section{The Rocq Theorem Prover}\label{coq_theorem}
The Rocq theorem prover~\cite{the_coq_development_team} is an interactive proof assistant that lets users define mathematical objects, state specifications about them, and prove that the definitions satisfy those specifications. Proof writing is the most challenging step; users guide Rocq using \emph{tactics}, a domain-specific language for building proofs incrementally. While Rocq offers some automation (e.g., decision procedures and proof search), human guidance is usually required to complete non-trivial proofs. Its popularity stems from dependent types, which allow encoding arbitrary mathematical statements directly in its logic, and from its ability to extract certified programs to efficient code in OCaml, Haskell, or Scheme.  Rocq has verified many real-world projects, including the CompCert C compiler~\cite{10.1145/1111037.1111042} (used by companies such as AbsInt\footnote{https://www.absint.com/partners.htm}), the Fiat-Crypto~\cite{8835346} library (used in BoringSSL for Chrome, Android, and CloudFlare), and the ConCert smart contract framework~\cite{10.1145/3372885.3373829}.

We use CertiRocq \cite{anand2017certicoq} that turns Rocq programs into WebAssembly (wasm) with machine-checked correctness guarantees. For wasm, CertiRocq extends its extraction pipeline to enable compilation of Rocq developments into \texttt{.wasm} binaries. The wasm backend, originally developed as CertiCoq-Wasm~\cite{meier2025certicoq} and now integrated into CertiRocq, is itself implemented and verified in Rocq with its semantics mechanised against the WebAssembly standard. It takes Rocq as input and compiles it to wasm code by mapping Rocq's primitive integer operations directly to native wasm instructions. The key insight is that programs proven correct within Rocq can be extracted, compiled, and executed in a wasm runtime while preserving their formal semantics. This bridges high-assurance formal verification with practical deployment, allowing certified components such as cryptographic protocols or voting primitives to run securely in web or sandboxed environments without sacrificing the guarantees established in the proof assistant. 




\section{Schnorr Protocol Rocq Encoding}\label{schn_proof}
  We briefly explain the Schnorr protocol, an interactive protocol 
  between a prover and a verifier. 
  Given some public input $(g, q, h)$ where $\langle g\rangle$
  is a cyclic group of prime order $q$ and $h\in\langle g\rangle$,
  the prover claims that they know a
  witness $w$ for the statement $h = g^w$; the existence of such
  a $w$ is immediate because $g$ generates the group. However,
  does the prover know the witness $w$? In order to convince the
  verifier, the prover and the verifier do the following:
  \begin{itemize}
    \item the prover picks a random number $u$, computes $a = g^u$,
    and sends $a$ to the verifier.
    \item the verifier picks a random challenge $c$ and sends it to
    the prover
    \item the prover computes $r = u + c \cdot w \pmod q$ and sends $r$ to the
    verifier
  \end{itemize}

The verifier accepts the transcript $(a, c, r)$ if $g^r = a \cdot h^c$, otherwise rejects. We can prove that 
this protocol satisfies completeness, special soundness, and special honest-verifier 
zero-knowledgeness. 


\begin{itemize}
  \item \textbf{completeness:} when $P$ and $V$ both follow the protocol, $V$ accepts the proof, i.e., 
  the verification equation checks out. We can see by algebraic simplification that it is the case.
    \begin{align*}
      g^r = g^{u + c \cdot w} = g^u \cdot (g^w)^c = a \cdot h^c
    \end{align*}
  \item \textbf{special soundness:} given two accepting 
  conversations $(a, c, r)$ and $(a, c', r')$ with $c \neq c'$,
  we have 
  \[ g^r = a \cdot h^c, \quad g^{r'} = a \cdot h^{c'}. \]
  It follows that $h = g^{(r-r')/(c-c')}$, hence we obtain
  the witness $w$ satisfying $h=g^w$ as $w = (r - r')/(c - c') \pmod q$.

  \item \textbf{special honest-verifier zero-knowledgeness:}
   given challenge $c$, the distribution of conversations between a prover using
   witness $w$ and an honest verifier and the distribution of conversations output by a
   simulator as defined below--without access to witness $w$--are identical. Both
   distributions (\ref{f}) and (\ref{snd}) are uniform with each conversation occurring 
   with probability $1/q$.
   \begin{align}\label{f}
    \{(a, c, r) \mid u \in_{R} Z_q ; a := g^u; r := u + c \cdot w \pmod q \}
   \end{align}
   \begin{align}\label{snd}
    \{(a, c, r) \mid r \in_{R} Z_q ; a := g^r \cdot h^{-c}\}
   \end{align}
   
\end{itemize}


In the Schnorr protocol, $g$, $h$, and $a$ are group elements from the underlying group, while $u$, $c$, and $r$ are scalar values from the underlying prime field. In concrete implementations, both the group elements and the field elements might be large 
numbers modulo some prime numbers. 
One possibility is to encode the Rocq functions to work with concrete numbers from the underlying group and field, but this leads to cumbersome proofs because we need to deal with the details of concrete numbers during the security proofs. More importantly, it is not modular because it binds the formalisation to specific groups, excluding the use elliptic curve groups, for instance. For most programming languages, this distinction 
may not be significant, but it can be a potential source of bugs \cite{10.1007/978-3-662-63958-0_24} undermining the security of underlying cryptographic primitives. 
However, when encoding it in Rocq, we need to thoroughly detail these distinctions because 
the same large number can belong to a group or field and that dictates what operations are available for them. 
Therefore, to avoid this unpleasant situation,  
we abstract the underlying group as an abstract type $G$, endowed with 
usual group operation, and the underlying field as an abstract type $F$, endowed with 
field operation in our formalisation. 
Moreover, we abstract these two types into a vector space for smooth interaction between 
group elements and field elements (exponentiation in a multiplicative group as well as multiplication of a field element 
with a point on an elliptic curve are instances of this interaction). This 
makes our formalisation generic and can be instantiated with 
any kind of group and field as long as it follows the axioms of 
a vector space. For example, in our formalisation we have developed 
an efficient algorithm for the Schnorr group but our formalisation can 
also be instantiated with an elliptic curve group. In addition, 
one of the major benefits of this abstraction is the ease with which proofs 
can be automated. Therefore, we model the Schnorr 
protocol, its simulator, and its verification equation in Rocq as shown in 
Listing \ref{schn_nat}. 
\begin{lstlisting}[frame=single, language=Coq, caption={Schnorr protocol},
  label={schn_nat},captionpos=t, basicstyle=\ttfamily\footnotesize,
  abovecaptionskip=-\medskipamount]
(*Real transcript (involves secret x)*)
Definition schnorr_protocol (x : F) (g : G) 
 (u c : F) : sigma_proto :=  
 ([g^u]; [c]; [u + c * x]).

(*Fake transcript (without the witness x)*)
Definition schnorr_simulator (g h : G) 
 (u c : F) : sigma_proto := 
 ([gop (g^u) (h^(opp c))]; [c]; [u]).

(* It checks if g^r = a * h^c *)
Definition accepting_conversation (g h : G) 
 (v : sigma_proto) : bool :=
 match v with
 | (a; c; r) =>  
    match Gdec (g^r) (gop a (h^c)) with 
    | left _ => true
    | right _ => false 
    end
 end.
\end{lstlisting}
  
\noindent
 $F$ is the underlying field, $G$ is the underlying group,
 $+ : F \rightarrow F \rightarrow F$ is the field addition,
 $* : F \rightarrow F \rightarrow F$ is the field multiplication, 
 $opp : F \rightarrow F$ is the additive inverse (field negation),
 $gop : G \rightarrow G \rightarrow G$ is the group operation, and
 $\mbox{\textasciicircum}: G \rightarrow F \rightarrow G$ is the scalar multiplication of the underlying vector space.  
Thus, in our setting we instantiate vectors with elements of $G$ and scalars with elements of $F$. 
Moreover, our notation for scalar multiplication is $\mbox{\textasciicircum}$ because it corresponds to exponentiation, 
when instantiated concretely with multiplicative groups, such as Schnorr groups.

We prove completeness and special soundness for the Schnorr protocol 
as shown in Listing \ref{comp_sound}. However, to prove the special honest-verifier 
zero-knowledge property we need a model to reason about uniform probability 
distribution in Rocq that we introduce in the next section. The source code for this section can be accessed from the
\href{https://anonymous.4open.science/r/SigmaProtocol-EBA5/src/Crypto/Sigma.v}{\textcolor{blue}{Sigma.v}} file from our anonymous repo.
 


\begin{lstlisting}[frame=single, language=Coq, caption={Completeness and Soundness},
  label={comp_sound},captionpos=t, basicstyle=\ttfamily\footnotesize,
  abovecaptionskip=-\medskipamount]
Lemma schnorr_completeness : 
forall (a : t G 1) (c r : t F 1),
(a; c; r) = (schnorr_protocol x g r c) ->
accepting_conversation g h (a; c; r) = true.
Proof. (* proof terms omitted *) Qed.
  
Lemma simulator_completeness : 
forall (a : t G 1) (c r : t F 1),
(a; c; r) = (schnorr_simulator g h r c) ->
accepting_conversation g h (a; c; r) = true.
Proof using -(x R). (* terms omitted *) Qed. 
  
Lemma special_soundness: 
forall (a : G) (ca ra cb rb : F), ca <> cb ->
accepting_conversation g h ([a]; [ca]; [ra]) ->  
accepting_conversation g h ([a]; [cb]; [rb]) ->
exists y : F, g^y = h /\ y = ((ra - rb) * inv (ca - cb)).
Proof using -(x R). (* proof terms omitted *) Qed.
  
\end{lstlisting}
 
\subsection{Uniform Distribution}
To reason about special honest-verifier zero-knowledge proof 
intuitively, we need a model for uniform probability distributions. 
We follow \cite{erwig2006functional}  and 
encode a probability distribution as shown in Listing \ref{prob_def}, using a list of pairs where the 
second element is the probability---represented as 
a rational number---for the first element.
Moreover, we define equality of distributions by 
stating that two distributions are equal if they are permutations 
of each other. This encoding is flexible and allows us 
to encode any kind of distribution. We prove that \texttt{dist} is a 
monad \cite{erwig2006functional} by 
giving it suitable definitions of \texttt{Ret} and 
\texttt{Bind}. Making it a monad allows us 
to compose a distribution $n$ times to 
obtain a distribution on a vector of $n$ elements, 
needed in proving the special honest-verifier 
zero-knowledge proof of $n$-fold compositions of $\Sigma$-protocols.
However, to reason about zero-knowledge properties of $\Sigma$-protocol, we need uniform 
distributions, and therefore we define a function 
\texttt{\seqsplit{uniform\_with\_replacement}} to generate uniform distributions as shown in Listing \ref{prob_def}.
The relevant source code can be accessed from \href{https://anonymous.4open.science/r/SigmaProtocol-EBA5/src/Probability/Prob.v}{\textcolor{blue}{Prob.v}} and \href{https://anonymous.4open.science/r/SigmaProtocol-EBA5/src/Probability/Distr.v}{\textcolor{blue}{Distr.v}}.

\begin{lstlisting}[frame=single, language=Coq, caption={Uniform Probability Distribution},
    label={prob_def},captionpos=t, basicstyle=\ttfamily\footnotesize,
    abovecaptionskip=-\medskipamount]
(* Probability Distribution on a type A *)
Definition dist (A : Type) : Type := list (A * prob).
  
(* Probability Monad *)
Definition Ret {A : Type} (x : A) : dist A := [(x, one)].

Fixpoint Bind {A B : Type} (xs : dist A)  
  (f : A -> dist B) : dist B := 
  match xs with 
  | [] => [] 
  | (ax, px) :: tx => 
    List.append (List.map (fun '(ut, pt) => 
    (ut, mul_prob px pt)) (f ax)) (Bind tx f)
  end.

(* Uniform distribution *)
Definition uniform_with_replacement {A : Type} : 
  forall (l : list A), l <> [] -> dist A.
Proof.
  intros ? H.
  remember (Pos.of_nat (List.length l)) as len.
  exact (List.map 
    (fun x => (x, mk_prob 1 len)) l).
Defined.
\end{lstlisting}

  
We now have all ingredients for the special honest-verifier zero-knowledge proof of the Schnorr protocol. We define the 
Schnorr distribution \texttt{\seqsplit{schnorr\_distribution}}
that involves the secret $x$ and 
the simulator distribution \texttt{\seqsplit{simulator\_distribution}}
that does not involve the secret $x$, see Listing~\ref{def_zero}. 
Both distributions sample a random field element $u$ from 
\texttt{\seqsplit{uniform\_with\_replacement}} \texttt{lf} \ \texttt{Hlfn}, 
where \texttt{lf} lists the elements from which $u$ is drawn and 
\texttt{Hlfn} asserts that \texttt{lf} is non-empty, 
to construct their respective distributions for a $\Sigma$-protocol.
Using these definitions, we prove that the two distributions coincide in the 
\texttt{\seqsplit{special\_honest\_verifier\_zkp}} theorem. Unlike the standard informal 
presentation pervasive throughout the cryptography literature, these distributions 
are not syntactically identical but they coincide modulo 
\texttt{\seqsplit{accepting\_conversation}}. We therefore apply \texttt{\seqsplit{accepting\_conversation}} 
to the first component (transcript) of each pair in both distributions, obtaining 
the boolean value \texttt{true} in all cases. After this transformation, the 
resulting distributions are identical.


\begin{lstlisting}[frame=single, language=Coq, caption={Special Honest-Verifier Zero-Knowledge Proof},
  label={def_zero},captionpos=t, basicstyle=\ttfamily\footnotesize,
  abovecaptionskip=-\medskipamount]
(* Distribution that involves the secret x *)
Definition schnorr_distribution  (lf : list F) 
  (Hlfn : lf <> List.nil) (x : F) (g : G) (c : F) : 
  dist sigma_proto :=
  (* draw u from a random distribution *)
  u <- (uniform_with_replacement lf Hlfn) ;;
  Ret (schnorr_protocol x g u c).

(* without secret x *)
Definition simulator_distribution (lf : list F) 
  (Hlfn : lf <> List.nil) (g h : G) (c : F) : 
  dist sigma_proto :=
  (* draw u from a random distribution *)
  u <- (uniform_with_replacement lf Hlfn) ;;
  Ret (schnorr_simulator g h u c).
  
Lemma special_honest_verifier_zkp : 
  forall (lf : list F) (Hlfn : lf <> List.nil) (c : F), 
  List.map (fun '(a, p) => 
    (accepting_conversation g h a, p))
    (@schnorr_distribution lf Hlfn x g c) = 
  List.map (fun '(a, p) => 
    (accepting_conversation g h a, p))
    (@simulator_distribution lf Hlfn g h c).
Proof. (* proof terms omitted *) Qed. 
\end{lstlisting}

As we stated previously, we have formalised methods for combining $n$ And, Or, Eq, Parallel, and Neq Schnorr statements, where proving the special honest-verifier zero-knowledge property is more challenging than for a single Schnorr statement. One major challenge is drawing $n$ random elements from a uniform distribution. We therefore define the function \texttt{\seqsplit{repeat\_dist\_ntimes\_vector}}, shown in Listing~\ref{dist_comp}, which takes a distribution $d$ and a number $n$ and returns a distribution over vectors of length $n$, encoded as \texttt{dist}\,(\texttt{Vector.t}\,A\,n). Note that this is a generic function that can be instantiated with any distribution. In cryptography, however, we are primarily concerned with the uniform distribution; in the theorem \texttt{\seqsplit{uniform\_probability\_multidraw\_prob}}, we instantiate the distribution with \texttt{uniform\_with\_replacement} \texttt{lf} \texttt{Hlf} and prove that its probability is $1/|lf|^n$. We use this function consistently throughout the formalisation. As an example, we illustrate a case of the parallel composition in which $n$ statements are combined, requiring a vector of $n$ random field elements to establish the zero-knowledge property, as shown in \texttt{\seqsplit{parallel\_schnorr\_distribution}}.


\begin{lstlisting}[frame=single, language=Coq, caption={Composition of a Distribution},
label={dist_comp},captionpos=t, basicstyle=\ttfamily\footnotesize,
abovecaptionskip=-\medskipamount]
Fixpoint repeat_dist_ntimes_vector {A : Type} 
(d : dist A) (n : nat) : dist (Vector.t A n) := 
match n with 
| 0 => Ret (@Vector.nil A)
| S n' => 
  Bind d (fun u => 
    Bind (repeat_dist_ntimes_vector d n')
    (fun v => Ret (Vector.cons _ u _ v)))
end.

Lemma uniform_probability_multidraw_prob 
{A : Type} : forall n (lf : dist A) 
(a : Vector.t A n) b (Hlf : lf <> []), 
In (a, b) (repeat_dist_ntimes_vector 
  (uniform_with_replacement lf Hlf) n) ->
b = mk_prob 1 (Nat.pow (List.length lf) n).
Proof. (* terms omitted *) Qed.

Definition parallel_schnorr_distribution  
  {n : nat} (lf : list F) (Hlfn : lf <> List.nil) 
  (x : F) (g : G) (cs : Vector.t F n) : 
  dist (@sigma_proto n n n) :=
  (* draw n random elements *)
  us <- repeat_dist_ntimes_vector 
    (uniform_with_replacement lf Hlfn) n ;;
  Ret (construct_parallel_schnorr x g us cs).
\end{lstlisting}


\subsection{OR Composition}
OR-compositions of $\Sigma$-protocols allow a prover to demonstrate knowledge of one witness among $n$ possible discrete logarithm statements 
without revealing which specific statement they know. Formally, given public group elements $(g_1, h_1), (g_2, h_2), \ldots, (g_n, h_n)$, 
the prover convinces the verifier that they know a witness $x_k$ such that $h_k = g_k^{x_k}$ for some $k \in \{1, \ldots, n\}$, 
without revealing the index $k$ or the witness $x_k$. The key insight behind the Or composition, introduced by Cramer, Damg\aa rd, and Schoenmakers~\cite{cramer1994proofs}, 
is to allow the prover a controlled degree of freedom to \emph{cheat} on statements for which they do not know witnesses.  We implement this protocol generically for $n \geq 2$ statements. The prover, knowing the witness for the $k$-th relation (where $k$ is specified as a finite index of sort \texttt{Fin.t (2+n)}), first enters the commitment phase by choosing random values: $u_k$ for the real commitment, and for each simulated statement $i \neq k$, random values $u_i$ and $c_i$. The prover then computes the real commitment $t_k = g_k^{u_k}$ for statement $k$ and simulated commitments $t_i = g_i^{u_i} \cdot h_i^{-c_i}$ for $i \neq k$, sending the vector of commitments $(t_1, \ldots, t_n)$ to the verifier. In the challenge phase, the verifier responds with a random challenge $c$. Finally, during the response phase, the prover computes $c_k = c - \sum_{i \neq k} c_i$, along with $s_k = u_k + c_k \cdot x_k$ as the honest response for the real statement, while for $i \neq k$, the responses $s_i$ are set to the previously chosen $u_i$. The prover then sends both response vectors $(s_1, \ldots, s_n)$ and $(c_1, \ldots, c_n)$. The verifier accepts if: (i) $c = \sum_{i=1}^{n} c_i$, and (ii) for each $i$, $g_i^{s_i} = t_i \cdot h_i^{c_i}$.


Our Rocq formalisation of the OR-composition of $\Sigma$-protocols supports $n \geq 2$ statements with arbitrary generators for each statement, enabling proofs of knowledge for relations $g_i^{x_i} = h_i$. The function \texttt{\seqsplit{generalised\_construct\_or\_conversations\_schnorr}} constructs a proof for the relation
$\{ (x_1, x_2, \ldots, x_n) \mid h_1 = g_1^{x_1} \lor h_2 = g_2^{x_2} \lor \cdots \lor h_n = g_n^{x_n} \}$, 
where the prover knows the witness for the relation indicated by the index \texttt{rindex}. One advantage of this generic formalisation is that we can instantiate \( n \) with any finite number of statements at compile time, thereby uniformly capturing arbitrary-length disjunctions without modifying the underlying protocol definition.  Similarly, we define the simulator \texttt{\seqsplit{generalised\_construct\_or\_conversations\_simulator}}, which generates accepting transcripts for the disjunctive relation without access to any witness. Building on this, we construct two distributions: \texttt{\seqsplit{generalised\_or\_schnorr\_distribution}}, which depends on the secret witness \(x\), and \texttt{\seqsplit{generalised\_or\_simulator\_distribution}}, which is independent of any witness. These distributions are then used to establish the zero-knowledge property captured by \texttt{\seqsplit{generalised\_or\_special\_honest\_verifier\_zkp}}. The two distributions are identical and therefore the protocol is information-theoretically zero-knowledge. The formalisation proves three security properties: completeness (honest execution always verifies), special soundness (which recovers a witness from two accepting transcripts with the same commitments but different challenges), and special honest-verifier zero-knowledge (the distribution of real transcripts involving the witness is identical to the simulator distribution without the witness). The source code for this section can be accessed in the \href{https://anonymous.4open.science/r/SigmaProtocol-EBA5/src/Crypto/OrSigmaGen.v}{\textcolor{blue}{OrSigmaGen.v}} file.

\begin{lstlisting}[frame=single, language=Coq, caption={Or-$\Sigma$-Protocol},
label={or_proof},captionpos=t, basicstyle=\ttfamily\footnotesize,
abovecaptionskip=-\medskipamount]
Definition generalised_construct_or_conversations_schnorr {n : nat} 
 (rindex : Fin.t (2 + n)) (x : F) (gs hs : Vector.t G (2 + n))
 (uscs : Vector.t F ((2 + n) + (1 + n))) (c : F) : 
 @sigma_proto F G (2 + n) (1 + (2 + n)) (2 + n).
Proof. (* proof terms omitted. *) Defined.

Definition generalised_construct_or_conversations_simulator {n : nat} 
 (gs hs : Vector.t G (2 + n)) (uscs : Vector.t F ((2 + n) + (1 + n)))  
 (c : F) :  @sigma_proto F G (2 + n) (1 + (2 + n)) (2 + n).
Proof. (* proof terms omitted. *) Defined.

(* schnorr OR distribution *)
Definition generalised_or_schnorr_distribution  {n : nat} 
 (lf : list F) (Hlfn : lf <> List.nil) (rindex : Fin.t (2 + n)) 
 (x : F) (gs hs : Vector.t G (2 + n)) (c : F) : 
 dist (@sigma_proto F G (2 + n) (1 + (2 + n)) (2 + n)) := 
 uscs <- repeat_dist_ntimes_vector 
  (uniform_with_replacement lf Hlfn) ((2 + n) + (1 + n)) ;;
 Ret (generalised_construct_or_conversations_schnorr 
  rindex x gs hs uscs c).

(* simulator distribution *)
Definition generalised_or_simulator_distribution {n : nat} 
 (lf : list F) (Hlfn : lf <> List.nil) 
 (gs hs : Vector.t G (2 + n)) (c : F) : 
 dist (@sigma_proto F G (2 + n) (1 + (2 + n)) (2 + n)) :=
 (* draw ((m + (1 + n)) + (m + n)) random elements *)
 usrs <- repeat_dist_ntimes_vector 
  (uniform_with_replacement lf Hlfn) ((2 + n) + (1 + n)) ;;
 Ret (generalised_construct_or_conversations_simulator 
  gs hs usrs c).

(* Information theoretic proofs *)
Lemma generalised_or_special_honest_verifier_zkp {n : nat}
 (x : F) (* secret witness *)
 (gs hs : Vector.t G (2 + n)) (rindex : Fin.t (2 + n))
 (* Prover knows rindex relation *) 
 (R : Vector.nth hs rindex = Vector.nth gs rindex ^ x) : 
 forall (lf : list F) (Hlfn : lf <> List.nil) (c : F),
 List.map (fun '(a, p) => 
  (generalised_or_accepting_conversations gs hs a, p))
  (generalised_or_schnorr_distribution lf Hlfn rindex x gs hs c) = 
 List.map (fun '(a, p) => 
 (generalised_or_accepting_conversations gs hs a, p))
 (generalised_or_simulator_distribution lf Hlfn gs hs c).
Proof. (* terms omitted *) Qed.

\end{lstlisting}

\subsection{Okamoto Protocol}

Okamoto's protocol~\cite[Identification scheme 1]{okamoto1992provably} was originally introduced as a provably secure variation of Schnorr's identification protocol.
While Schnorr's protocol proves knowledge of a single exponent $x$ such that $h = g^x$, Okamoto's protocol proves knowledge of 
exponents $x_1, x_2$ satisfying $h = g_1^{x_1} \cdot g_2^{x_2}$, given fixed generators $g_1,g_2$. Okamoto's protocol and its generalization to multiple exponents and generators
(already considered in~\cite{10.1007/3-540-39118-5_13}) are commonly viewed as $\Sigma$-protocols, and are often used as building block in more advanced constructions.
We will use it, in particular, for NEQ-composition as explained in the next section.

The generalized Okamoto protocol follows the standard 3-move $\Sigma$-protocol structure. Given public generators $\mathbf{g} = (g_1, g_2, \ldots, g_n)$, public value $h$, and secret witnesses $\mathbf{x} = (x_1, x_2, \ldots, x_n)$ satisfying $h = \prod_{i=1}^{n} g_i^{x_i}$, the prover begins by selecting random values from the underlying field $\mathbf{u} = (u_1, u_2, \ldots, u_n)$ and computing announcement $t = \prod_{i=1}^{n} g_i^{u_i}$, which is sent to the verifier. Next, the verifier responds with a random challenge $c$ from the field. Finally, the prover computes responses $r_i = u_i + c \cdot x_i$ for each $i \in {1, \ldots, n}$ and sends the vector $\mathbf{r} = (r_1, r_2, \ldots, r_n)$ to the verifier for verification. The verifier accepts if and only if 
$\prod_{i=1}^{n} g_i^{r_i} = t \cdot h^c$. Completeness follows from the fact that:
\[
\prod_{i=1}^{n} g_i^{r_i} = \prod_{i=1}^{n} g_i^{u_i + c \cdot x_i} = 
\left(\prod_{i=1}^{n} g_i^{u_i}\right) \cdot \left(\prod_{i=1}^{n} g_i^{x_i}\right)^c = t \cdot h^c
\]

Our Rocq formalisation implements the generalised Okamoto protocol for $n \geq 2$ generators and can be accessed in the file \href{https://anonymous.4open.science/r/SigmaProtocol-EBA5/src/Crypto/Okamoto.v}{\textcolor{blue}{Okamoto.v}}. The key functions include: \texttt{\seqsplit{generalised\_okamoto\_commitment}} (computes $\prod_{i=1}^{n} g_i^{u_i}$), \texttt{\seqsplit{generalised\_okamoto\_response}} (computes responses $r_i = u_i + c \cdot x_i$), and 
\texttt{\seqsplit{generalised\_okamoto\_accepting\_conversation}} (verifies the product equality condition). We also provide specialised instances for the base case of exactly two generators ($n=2$), which is heavily used in the Neq composition. The formalisation proves all three security properties: (i) \textbf{completeness}: 
when both prover and verifier follow the protocol honestly with $h = \prod_{i=1}^{n} g_i^{x_i}$, 
the verification check always succeeds. The proof proceeds by induction on $n$, leveraging the 
vector space properties of scalar multiplication distributivity over field addition; (ii) \textbf{special soundness}: 
given two accepting transcripts $(t, c_1, \mathbf{r}_1)$ and $(t, c_2, \mathbf{r}_2)$ with 
the same commitment $t$ but different challenges $c_1 \neq c_2$, an extractor can compute 
witnesses $\mathbf{x} = (x_1, \ldots, x_n)$ where 
$x_i = (r_{1,i} - r_{2,i}) \cdot (c_1 - c_2)^{-1}$ for each $i$. The proof establishes that 
these extracted values satisfy $h = \prod_{i=1}^{n} g_i^{x_i}$; and (iii) \textbf{special honest-verifier zero-knowledge}: 
the real transcript distribution (involving secrets $\mathbf{x}$) is identical to the simulator 
distribution (without secrets). The simulator chooses random $\mathbf{u}$ and $c$, then computes 
the commitment as $t = \left(\prod_{i=1}^{n} g_i^{u_i}\right) \cdot h^{-c}$ with responses 
$\mathbf{r} = \mathbf{u}$. Both distributions produce accepting transcripts with probability 
$1/|\mathbb{F}|^n$ where $\mathbb{F}$ is the underlying field, proven using our uniform probability 
distribution framework.


Beyond the standard $\Sigma$-protocol properties, our formalisation includes witness indistinguishability. We prove that transcripts generated with different witness 
vectors $\mathbf{x}_1$ and $\mathbf{x}_2$ satisfying the same relation $h = \prod_{i=1}^{n} g_i^{x_{1,i}} = 
\prod_{i=1}^{n} g_i^{x_{2,i}}$ are information-theoretically indistinguishable. The proof \texttt{\seqsplit{generalised\_okamoto\_witness\_indistinguishable}}, shown in Listing \ref{wit_ind}, relies on two 
bijective transformation functions that map 
randomness between the two witness spaces, along with the homomorphic property showing that 
$\prod_{i=1}^{n} g_i^{u_i - v_i} = \left(\prod_{i=1}^{n} g_i^{u_i}\right) \cdot 
\left(\prod_{i=1}^{n} g_i^{-v_i}\right)$.

\begin{lstlisting}[frame=single, language=Coq, caption={Witness Indistinguishability},
label={wit_ind},captionpos=t, basicstyle=\ttfamily\footnotesize,
abovecaptionskip=-\medskipamount]
Theorem generalised_okamoto_witness_indistinguishable {n : nat} :
  forall (gs : Vector.t G (2 + n)) (h : G) (xs1 xs2 : Vector.t F (2 + n)),
  (* Both witnesses produce the same public key h *)
  h = generalised_okamoto_commitment gs xs1 ->
  h = generalised_okamoto_commitment gs xs2 ->
  (* For every conversation (a, c, rs) from xs1, there exists unique us' 
    such that the same conversation is produced by xs2 *)
  forall (us rs : Vector.t F (2 + n)) (a : G) (c : F),
  ([a]; [c]; rs) = generalised_okamoto_real_protocol xs1 gs h us c ->
  (* Then there exists unique us' producing the same transcript for xs2 *)
  exists! us' : Vector.t F (2 + n),
  ([a]; [c]; rs) = generalised_okamoto_real_protocol xs2 gs h us' c.
Proof. (* proof terms omitted *) Qed.
\end{lstlisting}

\subsection{NEQ Composition}
The Neq composition formalises the statement that a prover knows witnesses
$x_1, \ldots, x_{2+n}$ for public instances $(g_i, h_i)$ such that $h_i = g_i^{x_i}$ for every
$i$, and moreover these witnesses are pairwise distinct, i.e., $x_i \neq x_j$ for all distinct pairs $i \neq j$. 
In contrast to Or and And compositions, which combine existing statements by logical connectives, the Neq protocol has to encode a family of inequality constraints between witnesses.  In our development, this is done by combining an And proof of all discrete-logarithm relations with one Okamoto proof for every unordered pair of witnesses.

The protocol structure for $2+n$ statements (with $n \ge 0$, i.e., at least two statements) is inherently quadratic in complexity due to the pairwise nature of inequality constraints. For each pair of indices $(i, j)$ with $i < j$, the prover must demonstrate that $x_i \neq x_j$ without revealing the values. Our implementation works as follows: in the commitment phase, the prover generates randomness of size $(2 + n) + ((2 + n) \cdot (1 + n))$ for two components. The first $2+n$ random values $u_i$ are used for the AND component, creating commitments $t_i = g_i^{u_i}$ for each $i \in \{1,\dots,2+n\}$, thereby proving knowledge of all discrete logarithms. The remaining randomness, of size $(2+n)(1+n)$, is split into two halves of equal length and zipped together to provide two random field elements per unordered pair $(i,j)$. For each pair, the prover generates an Okamoto commitment using the combined generators $(g_i \cdot g_j,\; h_i \cdot h_j)$ and the two allocated randomness values. 
The number of such commitments is $\binom{2+n}{2}$, one for each  pair $(i,j)$. The combined commitment vector, of dimension $(2+n) + \binom{2+n}{2}$, is then sent to the verifier. In the challenge phase, the verifier responds with a random challenge $c$ from the underlying field. In the response phase, the prover computes two sets of responses. The AND responses are $r_i = u_i + c \cdot x_i$ for each $i$ (as in the standard AND composition). For the Okamoto component, the prover constructs, for each pair $(i,j)$, the witnesses $w_{ij} = x_i \cdot (x_i - x_j)^{-1}$ and $w_{ji} = (x_j - x_i)^{-1}$ that encode the inequality $x_i \neq x_j$. Using these witnesses, the prover computes a two-dimensional Okamoto response; this yields $(2+n)(1+n)$ field elements. The total response vector thus has dimension $(2+n) + (2+n)(1+n) = (2+n)^2$, which is sent to the verifier.

The verifier accepts if both sub-protocol verifications succeed: (i) the AND verifier confirms that each $g_i^{r_i} = t_i \cdot h_i^c$, and (ii) for each pair $(i, j)$, the Okamoto verifier confirms that the relation $(g_i \cdot g_j)^{r_{i,1}} \cdot (h_i \cdot h_j)^{r_{i,2}} = t_{i,j} \cdot g_j^c$ holds, where $r_{i,1}, r_{i,2}$ are the two responses for pair $(i,j)$ and $t_{i,j}$ is the corresponding Okamoto commitment. The correctness of the Okamoto component for inequality relies on the algebraic identity: if $x_i = x_j$, then $(x_i - x_j)^{-1}$ would be undefined (division by zero), making it impossible for the prover to construct valid witnesses. Conversely, when $x_i \neq x_j$, the witnesses $\left(x_i \cdot (x_i - x_j)^{-1}, (x_j - x_i)^{-1}\right)$ satisfy $(g_i \cdot g_j)^{x_i \cdot (x_i - x_j)^{-1}} \cdot (h_i \cdot h_j)^{(x_j - x_i)^{-1}} = g_j$, where we use $h_i = g_i^{x_i}$, $h_j = g_j^{x_j}$, and $(x_j - x_i)^{-1} = -(x_i - x_j)^{-1}$.


Our Rocq formalisation, can be accessed in the file \href{https://anonymous.4open.science/r/SigmaProtocol-EBA5/src/Crypto/NeqSigma.v}{\textcolor{blue}{NeqSigma.v}}, proves the three standard $\Sigma$-protocol properties: (i) \textbf{completeness}: 
when the prover knows distinct witnesses $x_1, \ldots, x_{2+n}$ satisfying $g_i^{x_i} = h_i$ and $x_i \neq x_j$ 
for all $i \neq j$, the verification always succeeds. The proof combines completeness results from both the And composition completeness and the Okamoto protocol completeness for each pair; (ii) \textbf{special honest-verifier zero-knowledge}: 
the real transcript distribution (generated with witnesses $x_1, \ldots, x_{2+n}$) is identical to the 
simulator distribution (without witnesses). The simulator splits randomness between And and Okamoto 
components: for And, it uses the standard simulator producing commitments $t_i = g^{u_i} \cdot h_i^{-c}$; 
for each Okamoto pair, it generates simulated commitments independently. Both distributions produce 
accepting transcripts with probability $1/|\mathbb{F}|^{(2+n) + (2+n)(1+n)} = 1/|\mathbb{F}|^{(2+n)^2}$, proven 
using our uniform probability distribution framework; and (iii) \textbf{special soundness:} from two accepting transcripts sharing identical commitments but responding to distinct challenges \(c_1 \neq c_2\), an extractor is able to reconstruct witnesses \(x_1, \ldots, x_{2+n}\) satisfying the underlying discrete logarithm relations.  However, establishing that these extracted witnesses are genuinely distinct is far more subtle, and cannot, in general, be achieved without additional structural assumptions on the generators. In classical treatments, this is covered by an independence assumption; in our development, we instead expose the underlying algebraic relation explicitly. As formalised in Listing~\ref{neq_special_sound}, we obtain an unconditional dichotomy: either the extracted witnesses are pairwise distinct, or the generator structure itself is degenerate in a precise sense, namely that there exists a non-trivial discrete logarithm relation between generators, i.e., \(g_j = g_i^\alpha\) for some \(\alpha\) that is efficiently computable, or one of the generators collapses to the identity.  In this way, special soundness is derived without postulating independence as an axiom. Instead, we exploit decidability of the relevant algebraic predicates to transform what is typically an implicit assumption into an explicit disjunction:
the protocol is sound either because witnesses separate cleanly, or because the algebraic structure itself reveals an underlying dependency. We would like to emphasise the importance of Rocq in this development. For the special soundness proof, our initial expectation was simply to extract the witness vector from two accepting transcripts sharing identical commitments. However, during the proof, Rocq exposed additional structural requirements that were not apparent at the informal level, forcing us to confront hidden dependencies and edge cases that would otherwise have been silently assumed away.


The implementation \texttt{\seqsplit{generalised\_construct\_neq\_conversations\_real\_transcript}} requires randomness of dimension \((2+n) + (2+n)(1+n) = (2+n)^2\) and produces proof transcripts of size \(O(n^2)\), reflecting the inherently quadratic nature of enforcing pairwise inequality constraints. Consequently, while theoretically sound and fully general, this construction is not particularly practical for large-scale instances, as both the randomness requirements and transcript size grow rapidly with \(n\).

\begin{lstlisting}[frame=single, language=Coq, caption={Neq $\Sigma$ Protocol},
label={neq_special_sound},captionpos=t, basicstyle=\ttfamily\footnotesize,
abovecaptionskip=-\medskipamount]
(* size of proofs is O(n^2) for NEQ if we have n statements. *) 
Definition generalised_construct_neq_conversations_real_transcript 
 {n : nat} : Vector.t F (2 + n) -> Vector.t G (2 + n) -> 
 Vector.t G (2 + n) -> Vector.t F ((2 + n) + ((2 + n) * (1 + n))) 
 -> F -> @sigma_proto F G ((2 + n) + Nat.div ((2 + n) * (1 + n)) 2) 
  1 ((2 + n) + (2 + n) * (1 + n)).
Proof. intros xs gs hs us c. (* terms omitted *) Defined. 

Theorem generalised_neq_accepting_conversations_soundness :
  forall (n : nat) (a : Vector.t G (2 + n + (2 + n) * (1 + n) / 2)) 
  (c1 c2 : F) (rs1 rs2 : Vector.t F (2 + n + (2 + n) * (1 + n))) 
  (gs hs : Vector.t G (2 + n)), [c1] <> [c2] ->
  neq_accepting_conversations gs hs (a; [c1]; rs1) = true ->
  neq_accepting_conversations gs hs (a; [c2]; rs2) = true ->
  exists (xs : Vector.t F (2 + n)), 
  (forall (i : Fin.t (2 + n)), hs[@i] = gs[@i] ^ xs[@i]) /\
  ((forall (i j : Fin.t (2 + n)), i <> j -> xs[@i] <> xs[@j]) \/
  (exists (i j : Fin.t (2 + n)), i <> j /\ 
  ((exists (alpha : F), gs[@j] = gs[@i] ^ alpha) \/ 
  gs[@i] = gid \/ gs[@j] = gid)))
\end{lstlisting}

\subsection{Linear Relations}
A further type of subprotocol that we have formalised involves a simple linear relation between committed secret values.
Concretely, we consider the use of Pedersen commitments~\cite{pedersen1991non} for a standard setup given by two generators $g$ 
and $h$ (where the discrete logarithm of $h$ with respect to $g$ is unknown). 

For publicly given Pedersen commitments $C_1, C_2, \ldots, C_n$, publicly given coefficients $\alpha_1, \alpha_2, \ldots, \alpha_n$,
and public value $z$, the prover proves knowledge of values $(v_1, r_1), (v_2, r_2), \ldots, (v_n, r_n)$ such that:
\begin{itemize}
\item $C_i = g^{v_i} \cdot h^{r_i}$ for all $i$ (each commitment opens correctly)
\item $z = \sum_{i=1}^{n} \alpha_i \cdot v_i$ (a linear constraint on the committed values)
\end{itemize}

The protocol follows the standard three-move $\Sigma$-protocol structure for $n$ commitments with a single linear constraint, requiring $2n$ fresh randomness values (denoted $u_1, \ldots, u_n$ and $w_1, \ldots, w_n$) and producing $(n+1)$ commitment group elements and $2n$ response field elements. In the commitment phase, the prover samples fresh randomness $u_1, \ldots, u_n, w_1, \ldots, w_n$ and computes for each $i \in {1, \ldots, n}$ the commitment $t_i = g^{u_i} \cdot h^{w_i}$, mimicking the structure of $C_i$, along with an additional commitment $t_{n+1} = g^{\sum_{i=1}^{n} \alpha_i \cdot u_i}$ to the linear combination of randomness. The prover then sends the commitment vector $(t_1, \ldots, t_n, t_{n+1})$ to the verifier. In the challenge phase, the verifier responds with a random field element $c$. During the response phase, the prover computes for each $i \in \{1, 2, \dots, n\}$ the responses $s_{1,i} = u_i + c \cdot v_i$ and $s_{2,i} = w_i + c \cdot r_i$, sending the two vectors $(s_{1,1}, \ldots, s_{1,n})$ and $(s_{2,1}, \ldots, s_{2,n})$. The verifier accepts the transcript if both verification equations hold: for all $i$, $g^{s_{1,i}} \cdot h^{s_{2,i}} = t_i \cdot C_i^c$ (verifying each Pedersen commitment), and $g^{\sum_{i=1}^{n} \alpha_i \cdot s_{1,i}} = t_{n+1} \cdot g^{c \cdot z}$ (verifying the linear constraint). Verification correctness follows from the homomorphic properties of the commitment scheme: if the 
prover is honest with secrets $(v_i, r_i)$ satisfying $C_i = g^{v_i} \cdot h^{r_i}$ and 
$z = \sum_i \alpha_i \cdot v_i$, then:
\begin{align*}
g^{s_{1,i}} \cdot h^{s_{2,i}} &= g^{u_i + c \cdot v_i} \cdot h^{w_i + c \cdot r_i} = 
(g^{u_i} \cdot h^{w_i}) \cdot (g^{v_i} \cdot h^{r_i})^c = t_i \cdot C_i^c \\
g^{\sum_i \alpha_i \cdot s_{1,i}} &= g^{\sum_i \alpha_i \cdot (u_i + c \cdot v_i)} = 
g^{\sum_i \alpha_i \cdot u_i} \cdot g^{c \cdot \sum_i \alpha_i \cdot v_i} = t_{n+1} \cdot g^{c \cdot z}
\end{align*}

Our Rocq formalisation, available at \href{https://anonymous.4open.science/r/SigmaProtocol-EBA5/src/Crypto/PedLinearRel.v}{\textcolor{blue}{PedLinearRel.v}}, implements this protocol with generic support for arbitrary $n \geq 2$ commitments and linear constraints. The formalisation proves all three security properties: (i) \textbf{completeness}: when the prover knows values $(v_1, r_1), \ldots, (v_n, r_n)$ satisfying both $C_i = g^{v_i} \cdot h^{r_i}$ and $z = \sum_i \alpha_i \cdot v_i$, the protocol always verifies successfully; (ii) \textbf{special soundness}: from two accepting transcripts with the same commitments but different challenges $c_1 \neq c_2$, 
an extractor recovers the witness values: $v_i = (s_{1,i}^{(1)} - s_{1,i}^{(2)}) \cdot (c_1 - c_2)^{-1}$ 
and similarly for $r_i$. The extraction is guaranteed to produce valid commitments $C_i = g^{v_i} \cdot h^{r_i}$. 
For the linear constraint, the formalisation proves a dichotomy: either $z = \sum_i \alpha_i \cdot v_i$ holds, 
or the group element $g$ is the identity (which is a degenerate case). This unconditional result avoids 
assumptions on the problem structure; and (iii) \textbf{special honest-verifier zero-knowledge}: 
the distribution of real transcripts (generated with honest witnesses) is identically equal to the 
simulator distribution (without witnesses). Both produce accepting transcripts with probability 
$1/|\mathbb{F}|^{2n}$, proven using our uniform distribution framework. The identical distributions 
guarantee information-theoretic security: the verifier learns nothing beyond the validity of the 
statements, even with unlimited computational power.

The above $\Sigma$-protocol is particularly powerful in blockchain privacy applications, such as confidential transactions~\cite{10.1007/978-3-030-17653-2_22}. In systems like Bitcoin, transaction amounts are publicly visible, which compromises user privacy. To address this, a sender (e.g., Alice) can commit to the input amount $v_1$, the output amount $v_2$, and her new balance $v_3$ using Pedersen commitments of the form $C_i = g^{v_i} h^{r_i}$. She then proves that the committed values satisfy the linear relation $v_1 - v_2 - v_3 = 0$ (i.e., money is conserved) without revealing any individual amount. Using coefficients $\alpha = (1, -1, -1)$ and public value $z = 0$, the proof establishes that the input balance equals the sum of the transferred amount and the resulting balance. This guarantees that no money is created or destroyed, while all transaction values remain hidden. We formalise this idea in the \href{https://anonymous.4open.science/r/SigmaProtocol-EBA5/src/Examples/PedLinearRelIns.v}{\textcolor{blue}{PedLinearRelIns.v}} file. The same technique is employed in real-world cryptocurrencies such as Mimblewimble and Monero. The protocol supports arbitrary linear constraints through the algebraic structures provided by the field and group abstractions, making it widely applicable to systems requiring commit-and-prove mechanisms over structured secrets.

Note that sometimes the direct approach of just revealing $r=\sum_{i=1}^{n} \alpha_i \cdot r_i$ and checking that $g^z h^r = \prod_{i=1}^n C_i^{\alpha_i}$ suffices in standalone protocols (see, e.g, the tallying protocol of~\cite{CFSY96}).

\subsection{Efficient Modular Arithmetic}
As stated earlier, to ensure the ease of proving and modularity, 
we worked with an (abstract) vector space. However, to execute the code of 
our formalisation, we need to instantiate the vector space with a concrete entity. 
Therefore, we have encoded the Schnorr group \cite{schnorr1991efficient} (multiplicative 
group of integers modulo a large prime p), shown in Listing \ref{group_def}, and a prime field, shown in Listing \ref{field_def}, in the Rocq theorem prover. 
To encode the Schnorr group, we postulate three integers $k$, $p$, and $q$. 
We assume $p$ and $q$ are primes and $p = k q + 1$ (when 
$k = 2$, $p$ is called a safe prime and $q$ is a Sophie Germain prime). Finally, we encode the neutral element (\texttt{one}), the group multiplication (\texttt{\seqsplit{mul\_schnorr\_group}}), 
and the group inverse (\texttt{\seqsplit{inv\_schnorr\_group}}) and we prove that our encoding 
of the Schnorr group in Rocq is a commutative group by making it 
an instance of the typeclass \texttt{\seqsplit{schnorr\_comm}}, which is an abstract typeclass that bundles the axioms of a commutative group.


 \begin{lstlisting}[frame=single, language=Coq, caption={Rocq Encoding of the Schnorr Group},
  label={group_def},captionpos=t, basicstyle=\ttfamily\footnotesize,
  abovecaptionskip=-\medskipamount]
Section SchnorrGroup. 
  Context {k p q : Z} {Hk : p = k * q + 1}
    {Hp : prime p} {Hq : prime q}.

  Record Schnorr_group := mk_schnorr {v : Z; Ha : 0 < v < p; 
    Hb : v ^ q mod p = 1}.

  (* u ^ (p - 2) is inverse of u *)
  Definition inv_schnorr_group (u : Schnorr_group) : Schnorr_group.
  refine(match u with 
    | mk_schnorr au Hu Hv => mk_schnorr (* terms omitted *)
    end).
  Defined.

  (* Schnorr is a commutative Group *)
  Global Instance schnorr_comm : 
    @commutative_group Schnorr_group (@eq Schnorr_group) 
    mul_schnorr_group one inv_schnorr_group.
End SchnorrGroup.
\end{lstlisting}


In order to encode a prime field, we postulate an integer $p$ and 
axiom that $p$ is prime. We encode the prime field as 
a record \texttt{Zp} that contains an integer $v$ with an axiom 
that it is in the range $[0 \ldots p)$. Moreover, 
we define all the operations of \texttt{Zp} and prove that 
it is an instance of the typeclass \texttt{zp\_field}, which is 
a typeclass that encodes the axioms of a field.

\begin{lstlisting}[frame=single, language=Coq, caption={Rocq Encoding of Prime Field},
  label={field_def},captionpos=t, basicstyle=\ttfamily\footnotesize,
  abovecaptionskip=-\medskipamount]
(* Prime Field *)
Section ZpField.
  Context {p : Z} {Hp : prime p}.

  Record Zp := mk_zp {v : Z; Hv : Z.modulo v p = v}.

  Global Instance zp_field : @field Zp zero 
    one zp_opp zp_add zp_sub zp_mul zp_inv zp_div.
End ZpField. 
\end{lstlisting}

Finally, we combine the Schnorr group and the prime field 
into a vector space instantiating the vector set with 
the Schnorr group and the scalar set with the prime field, 
shown in Listing \ref{vec_space}. 
However, in order to establish that it is an 
instance of vector space we need the definition 
of scalar multiplication. We encode the  
the scalar multiplication \texttt{pow} of 
vector space that takes an element from 
a Schnorr group and an element from 
the prime field and returns an element in the 
Schnorr group.  The reader can see that our encoding eliminates all 
possibilities of mixing group and field, both represented 
as an integer modulo some large prime \cite{10.1007/978-3-662-63958-0_24}.
The source code for this section can be accessed in the file \href{https://anonymous.4open.science/r/SigmaProtocol-EBA5/src/Utility/Zpstar.v}{\textcolor{blue}{Zpstar.v}}.

\begin{lstlisting}[frame=single, language=Coq, caption={Vector Space Instantiation},
  label={vec_space},captionpos=t, basicstyle=\ttfamily\footnotesize,
  abovecaptionskip=-\medskipamount]

  Definition pow (g : @Schnorr.Schnorr_group p q) 
  (y : @Zpfield.Zp q) : @Schnorr.Schnorr_group p q.
  refine 
    match g, y with 
    | mk_schnorr gt Hgta Hgtb, mk_zp yt Hyt => 
      mk_schnorr  (Npow_mod gt yt p) 
      (pow_subproof_first _ _ Hgta Hyt)
      (pow_subproof_second _ _ Hyt Hgta Hgtb)
    end.
  Defined.

  Global Instance pow_vspace : 
  @vector_space 
    (* Field *)
    (@Zpfield.Zp q) (@eq (@Zpfield.Zp q))
    (@Zpfield.zero q Hq) (@Zpfield.one q Hq)
    (@Zpfield.zp_add q) (@Zpfield.zp_mul q)
    (@Zpfield.zp_sub q) (@Zpfield.zp_div q)
    (@Zpfield.zp_opp q Hq) (@Zpfield.zp_inv q)
    (* Vector *)
    (Schnorr_group p q) (@eq (Schnorr_group p q))
    (@one p q Hp Hq) 
    (@inv_schnorr_group k p q Hk Hp Hq)
    (@Schnorr.mul_schnorr_group p q Hp Hq)
    pow.
  
\end{lstlisting}

\subsection{Efficient Extractor}
In this section, we explain the key ideas underlying the formalisation of an efficient extractor and simulator. Recall that for a $\Sigma$-protocol, the extractor is required to be efficient, i.e., it runs in time polynomial in the size of its input. However, functions defined over vector spaces do not capture computational complexity; in particular, they provide no information about the number of computational steps required to evaluate them on a given input. One possible approach in our formalisation would have been to adopt a monadic style~\cite{10.1145/3473585,10.1145/3158124}, where each function additionally returns the number of steps taken during its execution. This approach has the advantage of compositionality: given two monadic functions, their composition naturally yields the complexity of the combined computation. However, during code extraction, this additional bookkeeping would propagate through all functions, significantly complicating and effectively degrading the intended API. Instead, we adopt a more targeted approach. In practice, an extractor over a multiplicative group relies on the ability to efficiently compute inverses. In our concrete model, inverse is implemented via binary exponentiation using the function \texttt{\seqsplit{repeat\_op\_ntimes\_rec}}. We therefore introduce an inductive data type \texttt{\seqsplit{complexity\_repeat\_op\_ntimes\_rec}}, shown in Listing~\ref{extractor_complex}, that captures the number of steps performed by this binary exponentiation procedure. Finally, we prove in theorem \texttt{\seqsplit{correct\_complexity\_repeat\_op\_ntimes\_rec}} that the function \texttt{\seqsplit{repeat\_op\_ntimes\_rec}} runs in logarithmic time in the value of its input, equivalently linear in the bit-length of the input. This approach allows us to preserve a clean API without introducing monadic overhead. While it lacks the compositional advantages of the monadic approach, it is sufficient for our purposes, as our primary objective is to extract executable (OCaml and WebAssembly) code whose APIs closely match those of real-world cryptographic libraries. The source code for this section can be access in the \href{https://anonymous.4open.science/r/SigmaProtocol-EBA5/src/Utility/Functions.v}{\textcolor{blue}{Functions.v}} file.


\begin{lstlisting}[frame=single, language=Coq, caption={Extractor Complexity},
  label={extractor_complex},captionpos=t, basicstyle=\ttfamily\footnotesize,
  abovecaptionskip=-\medskipamount]
Inductive complexity_repeat_op_ntimes_rec 
(e t : N) : positive -> N -> N -> Prop :=
| xH_case : complexity_repeat_op_ntimes_rec e t xH 
    (N.modulo e t) N0
| xO_case p ret w : 
    complexity_repeat_op_ntimes_rec e t p ret w ->
    complexity_repeat_op_ntimes_rec e t (xO p) 
    (N.modulo (ret * ret) t) (N.succ w)
| xI_case p ret w : 
    complexity_repeat_op_ntimes_rec e t p ret w ->
    complexity_repeat_op_ntimes_rec e t (xI p) 
    (N.modulo (e * (ret * ret)) t) (N.succ w).

Lemma correct_complexity_repeat_op_ntimes_rec : 
  forall n e w, complexity_repeat_op_ntimes_rec e w n 
    (repeat_op_ntimes_rec e n w)  (N.log2 n).
Proof. (* proof terms omitted *) Qed.
\end{lstlisting}

\section{Case Studies}\label{case_studies}

\subsection{Approval Voting}
Approval voting is a voting method where each voter can approve any subset of candidates independently. 
Unlike plurality voting where each voter selects exactly one candidate, approval voting allows voters to express 
support for multiple candidates, resulting in fairer representation of voter preferences. Our formalisation 
implements approval voting using ElGamal encryptions paired with OR-compositions of $\Sigma$-proofs to guarantee that 
each vote is either 0 (not approved) or 1 (approved).  The frontend functions, shown in Listing \ref{frontend_fun}, realise the ballot-construction pipeline for approval voting. The function \texttt{\seqsplit{encrypt\_vote}} computes the ElGamal ciphertext $(g^r,\; g^m \cdot h^r)$ from generator $g$, public key $h$, vote value $m$, and randomness $r$. The function \texttt{\seqsplit{generate\_enc\_proof}} then constructs a two-branch $\Sigma$-proof that the ciphertext encrypts either zero ($0$) or a non-zero value: it performs a case split with \texttt{Fdec} $m$ \texttt{zero}, places the real witness in the corresponding branch index (\texttt{Fin.F1} when $m=0$, otherwise \texttt{Fin.FS} \texttt{Fin.F1}), and simulates the other branch. Finally, \texttt{\seqsplit{encrypt\_vote\_and\_generate\_enc\_proof}} composes both steps and returns the ciphertext together with its proof as a single object for verification. Ideally, for the non-zero case a voter should encrypt $1$ but it runs on the voter's computer so the voter can encrypt any value and produce a zero-knowledge proof but the verification function running at the backend (server) will reject the proof, except for some negligible probability. More precisely, theorem \texttt{\seqsplit{vote\_proof\_invalid\_reject}} shows that for invalid votes $m \neq 0,1$, verification returns false under the usual non-degenerate assumptions. The companion theorem \texttt{\seqsplit{vote\_proof\_invalid\_accept}} captures the exceptional case: if the prover can predict the verifier challenge ($c_1 = c$), an invalid transcript can still be accepted. Together, these results characterise the standard $\Sigma$-protocol soundness boundary and explain why unpredictable challenges are essential \cite{10.1007/978-3-642-34961-4_38,dao2023weak}. 

Finally, the function \texttt{\seqsplit{encrypt\_ballot\_and\_generate\_enc\_proof}} extends vote encryption to the entire ballot by vectorizing across all $n$ candidates. Each candidate 
produces a tuple $(C_i, \pi_i)$ where $C_i$ is the encrypted vote and $\pi_i$ is the accompanying $\Sigma$-proof. A complete ballot is thus a vector of $n$ such tuples. We extract all frontend functions into WebAssembly. The related Rocq proofs can be accessed from the \href{https://anonymous.4open.science/r/SigmaProtocol-EBA5/src/Frontend/Approval.v}{\textcolor{blue}{Approval.v}} file and the corresponding WebAssembly file from \href{https://anonymous.4open.science/r/SigmaProtocol-EBA5/src/Wasm/approval.wasm}{\textcolor{blue}{approval.wasm}}.


\begin{lstlisting}[frame=single, language=Coq, caption={Frontend Functions},
  label={frontend_fun},captionpos=t, basicstyle=\ttfamily\footnotesize,
  abovecaptionskip=-\medskipamount]
(* encrypt vote *)
Definition encrypt_vote (g h : G) (r : F) (m : F) : G * G :=
  @enc F G gop gpow g h m r.

(* create a proof *)
Definition generate_enc_proof 
  (g h : G) (r : F)  (m : F) (uscs : Vector.t F 3) (cp : G * G) 
  (c : F) : @Sigma.sigma_proto F (G * G) 2 3 2 :=
    match Fdec m zero with 
    | left _ => 
      @generalised_construct_encryption_proof_elgamal_real F zero 
        add mul sub opp G ginv gop gpow 0 Fin.F1 r uscs 
        [g^m; g^one] g h cp c (* real one goes to 
        the first place Fin.F1 *)
    | right _ =>  
      @generalised_construct_encryption_proof_elgamal_real F zero 
        add mul sub opp G ginv gop gpow 0 (Fin.FS Fin.F1) r uscs 
        [g^zero; g^m] g h cp c (* real one 
        goes to the second place (Fin.FS Fin.F1) *)
    end.

Definition encrypt_vote_and_generate_enc_proof 
  (g h : G) (r : F)  (m : F) (uscs : Vector.t F 3) (c : F) : 
  G * G * @Sigma.sigma_proto F (G * G) 2 3 2 := 
  let cp := @enc F G gop gpow g h m r in
  let sig_proof := generate_enc_proof g h r m uscs cp c
  in (cp, sig_proof).

Theorem vote_proof_invalid_reject : 
 forall (r : F) (g h : G) (m : F)  (u1 u2 c1 : F) (c : F),
 g <> gid -> (m <> zero /\ m <> one) -> c1 <> c ->
 verify_encryption_vote_proof g h 
 (encrypt_vote_and_generate_enc_proof g h r m [u1; u2; c1] c) = false
Proof. (* proof terms omitted. *) Qed. 

Theorem vote_proof_invalid_accept :
 forall (r : F) (g h : G) (m : F)  (u1 u2 c1 : F) (c : F),
 (m <> zero /\ m <> one) -> c1 = c ->
 verify_encryption_vote_proof g h 
 (encrypt_vote_and_generate_enc_proof g h r m [u1; u2; c1] c) = true.
Proof. (* proof terms omitted. *) Qed.
\end{lstlisting}


\textbf{Homomorphic Tallying}: Tallying is specified as an inductive state machine \texttt{count}, parameterised 
over another inductive datatype \texttt{state}, similar to \cite{10.1007/978-3-319-66107-0_26}. The \texttt{state} datatype has two constructors, \texttt{partial} and \texttt{finished}, that signifies that the counting of ballots is still ongoing or finished. More precisely, \texttt{partial} \texttt{us} \texttt{vbs} \texttt{inbs} \texttt{ms}: the counting is in progress, with all seen ballots \texttt{us}, valid ballots \texttt{vbs}, invalid ballots \texttt{inbs}, and running encrypted tally \texttt{ms}; and \texttt{finished} \texttt{us} \texttt{vbs} \texttt{inbs} \texttt{pt}: the computation is finalised and returns plaintext tally vector \texttt{pt}.

The judgment \texttt{count : state -> Type} defines legal counting traces through four rules: 
(i) \textbf{\texttt{ax}} (initialisation): starts from empty ballot lists and an identity tally vector \texttt{ms} $= (1_G,1_G),\ldots,(1_G,1_G)$. This means no vote has contributed yet; 
(ii) \textbf{\texttt{cvalid}} (accept valid ballot): if a new ballot \texttt{u} satisfies \texttt{\seqsplit{verify\_encryption\_ballot\_proof}} \texttt{u = true}, then \texttt{u} is moved into \texttt{vbs} (set of valid ballots), and the encrypted tally is updated homomorphically: \texttt{nms = \seqsplit{mul\_encrypted\_ballots}} \texttt{ms} (\texttt{map fst u}). Intuitively, ciphertext components are multiplied coordinate-wise, so exponents (hence votes) add; 
(iii) \textbf{\texttt{cinvalid}} (reject invalid ballot): if \texttt{\seqsplit{verify\_encryption\_ballot\_proof}} \texttt{u = false}, then \texttt{u} is appended to \texttt{inbs} (set of invalid ballots), while the running (encrypted) tally \texttt{ms} remains unchanged; and 
(iv) \textbf{\texttt{cfinish}} (finalisation): after all ballots are processed, decryption outputs \texttt{ds} and decryption proofs \texttt{pf} are checked. The rule requires pointwise relation $g^{pt_i}=ds_i$, i.e., plaintext counts are discrete logs of decrypted group elements and acceptance of the decryption proof \texttt{\seqsplit{decryption\_proof\_accepting\_conversations\_vector}} \texttt{ms} \texttt{ds} \texttt{pf = true}. Only then the system transitions to \texttt{finished}.

The executable function \texttt{\seqsplit{compute\_final\_tally}} implements the first phase (classification and encrypted aggregation) by recursion on the ballot list. For each ballot, it branches on \texttt{\seqsplit{verify\_encryption\_ballot\_proof}}: the true branch applies \texttt{cvalid}, the false branch applies \texttt{cinvalid}. It returns \texttt{(vbs,inbs,ms)} together with the invariant \texttt{Permutation bs (vbs ++ inbs)}, ensuring that every input ballot is accounted for exactly once. Finally, the function \texttt{\seqsplit{compute\_final\_count}} then derives the final plaintext tally by decrypting aggregate ciphertexts with secret key $x$ to obtain \texttt{ds}, computing \texttt{pt} by pointwise discrete-log search relative to generator $g$, and concludes \texttt{count (finished bs vbs inbs pt)} via \texttt{cfinish}. Hence, the formal counting pipeline is: \emph{verify each ballot} $\to$ \emph{homomorphically aggregate valid ciphertexts} $\to$ \emph{verify decryption} $\to$ \emph{extract plaintext counts}, with explicit proof obligations at each step. The source code for this section can be accessed in the \href{https://anonymous.4open.science/r/SigmaProtocol-EBA5/src/Backend/Tally.v}{\textcolor{blue}{Tally.v}} file.


\subsection{Helios Verifier}
To demonstrate the practical applicability of our formalisation, we have developed a certified verifier for 
the Helios voting system that processes real-world election data from the IACR elections. Helios is a 
widely-deployed open-source internet voting system that implements end-to-end verifiable elections using 
cryptographic proofs. Our verifier validates the complete election tally, including ballot verification, ciphertext aggregation, and decryption correctness using formally verified proofs 
extracted as OCaml code. 

The election verification process follows a very similar approach described in the 
previous section. We have an inductive datatype \texttt{count}, index by \texttt{state}. 
The \texttt{count} and \texttt{state} are similar to previously defined, except the fourth constructor \texttt{finish} is 
slightly different. The constructor \textbf{finish}: after all ballots are processed, the verifier performs end-to-end 
verification checks: (i) Decryption verification: Each of the $m$ talliers has published a decryption factor $d_i = \alpha^{s_i}$ 
(where $\alpha$ is the first component of the final ciphertext and $s_i$ is their secret key). The verifier 
checks the decryption factor's validity using Chaum-Pedersen proofs that prove the discrete log relation 
$\log_g(y_i) = \log_\alpha(d_i)$ (same exponent used for both); (ii) Plaintext verification: The plaintext tally is computed from decryption factors and checked for 
consistency, i.e., for each candidate position, the prover must verify that $g^{\text{plaintext}} = \text{decrypted value}$; (iii) Proof of knowledge (POK) checks: Each tallier has published a Schnorr proof of knowledge of their 
secret key, which is verified to ensure key ownership; and (iv) Public key consistency: The verifier confirms that the product of all talliers' public keys equals 
the election's public key $h = \prod_i y_i$, ensuring no key material was substituted.  The function \texttt{compute\_final\_count} is the top-level certified
verifier. Its type states that for a list of ballots $\mathit{bs}$,
a list of talliers $\mathit{ts}$, and a published plaintext
tally $\mathit{pt}$, one can always compute a Boolean $\mathit{bfinal}$ together
with a partition of $\mathit{bs}$ into valid ballots $\mathit{vbs}$ and invalid
ballots $\mathit{inbs}$, and---crucially---a term of type
$\texttt{count}~(\texttt{finished}~\mathit{bs}~\mathit{vbs}~\mathit{inbs}~\mathit{ts}~\mathit{pt}~\mathit{bfinal})$.
Because \texttt{count} is an inductive datatype whose constructors encode all
verification rules, inhabiting this type is itself a machine-checked
certificate of the entire end-to-end verification. The related Rocq proof can be accessed in the \href{https://anonymous.4open.science/r/SigmaProtocol-EBA5/src/Backend/HeliosTally.v}{\textcolor{blue}{HeliosTally.v}} file.

The generic function \texttt{compute\_final\_count} is parameterised over an
arbitrary field $F$, group $G$, and all associated operations.
The function \texttt{compute\_final\_count\_ins} fixes every one of these
parameters to the concrete Helios cryptographic setting: the field $F$ is
instantiated to $\mathbb{Z}_q$ (via the \texttt{Zpfield} module), the group
$G$ is the Schnorr group $\mathbb{Z}^*_p$ of order $q$ (via the
\texttt{Schnorr\_group} module), and the group operations---identity,
inverse, multiplication, and exponentiation---are the corresponding
modular-arithmetic functions \texttt{Schnorr.one}, \texttt{\seqsplit{inv\_schnorr\_group}},
\texttt{\seqsplit{mul\_schnorr\_group}}, and \texttt{pow}.
The generator $g$ and the election public key $h$
are supplied as further arguments, together with the ballot list $\mathit{bs}$,
the tallier list $\mathit{ts}$, and the published plaintext tally
$\mathit{pt}$.
The body of \texttt{\seqsplit{compute\_final\_count\_ins}} is a single call to
\texttt{\seqsplit{compute\_final\_count}} with all arguments filled in; the result is a
term of the fully instantiated type
$\texttt{count}~(\texttt{finished}~\mathit{bs}~\mathit{vbs}~\mathit{inbs}~\mathit{ts}~\mathit{pt}~\mathit{bfinal})$,
which is a machine-checked certificate that the election tally is correct. The related Rocq code can be accessed from  the 
\href{https://anonymous.4open.science/r/SigmaProtocol-EBA5/src/Examples/HeliosTallyIns.v}{\textcolor{blue}{HeliosTallyIns.v}} file.


Finally, the toplevel \texttt{main.ml} file ties together the parser
and the certified verifier (extracted from Rocq proofs) into a single end-to-end pipeline.
The program reads election data from standard input, which is structured
in three sections: the list of voter ballots
(\texttt{bs}), the list of trustees (\texttt{us}), and the published
plaintext tally (\texttt{rs}).
The parser reconstructs each ballot as a tuple containing the cast timestamp,
a vector of ElGamal ciphertexts $(\alpha_i, \beta_i)$, a vector of $\Sigma$-proofs (one per candidate, certifying that the plaintext is $0$ or $1$),
and metadata strings (election hash, voter hash, etc.).
Each trustee is parsed into a tuple containing the decryption factor vector,
the Chaum-Pedersen decryption proof vector, the Schnorr proof of knowledge,
and the trustee's public key $(g, p, q, y_i)$. Once parsing is complete, \texttt{compute\_final\_count\_ins} is invoked
with the concrete Helios parameters:
the number of candidates $n = 7$, the number of additional talliers
$m = 2$ (so that the total number of talliers is $1 + m = 3$), the
2024 election public key \texttt{h2024}, and the three parsed inputs
\texttt{bs}, \texttt{us}, \texttt{rs}.
The return value is a nested \texttt{existT}---the Curry-Howard
image of the $\Sigma$-type---which is pattern-matched to recover the
Boolean verdict \texttt{bfinal}, the lists of valid and invalid ballots
\texttt{vbs} and \texttt{inbs}, and, crucially, a \texttt{count} term
that is a machine-checked certificate of the entire verification.
The program then prints the full certificate trace by
recursively deconstructing the \texttt{count}, displaying for
each ballot whether it was valid or invalid and reporting the four individual Boolean checks---decryption proof
validity, plaintext consistency, proof-of-knowledge validity, and public-key product consistency.

The verifier has successfully validated real IACR 2024 and IACR 2023 election data, demonstrating practical deployment 
of formal verification in privacy-preserving voting systems. The end-to-end verification provides strong 
cryptographic assurance: if any ballot is modified, any decryption factor is forged, or any tally is 
changed, the verifier will detect the corruption. The related source code can accessed 
from the \href{https://anonymous.4open.science/r/SigmaProtocol-EBA5/src/Executable/HeliosVerifiercode/main.ml}{\textcolor{blue}{main.ml}} file.


\textbf{Primality of $p$ and $q$}:
The two large primes used in the Helios parameters are a 256-bit prime~$q$
and a 2048-bit prime~$p = kq + 1$.
Their primality is established inside Rocq using the
\textsc{Coqprime}~\cite{thery2007primality} library.  For $q$, the proof can be accessed from the \href{https://anonymous.4open.science/r/SigmaProtocol-EBA5/src/Examples/primeQ.v}{\textcolor{blue}{primeQ.v}} file and 
for $p$, the proof can be accessed from \href{https://anonymous.4open.science/r/SigmaProtocol-EBA5/src/Examples/primeP.v}{\textcolor{blue}{primeP.v}} file.  



  


\section{Related Work, Limitations, and Future Work}\label{rel_work}
Given the widespread use of $\Sigma$-protocols in electronic voting and related domains, their extensive formalisation in prior work is unsurprising~\cite{5552642,9519460,butler2021formalising,10.1145/3319535.3354247,287095,almeida2010certifying,10221929,10.1145/3594735}.  Barthe et al.~\cite{5552642} present the first formalisation of $\Sigma$-protocols and their composition in CertiCrypt, a library developed in the Rocq theorem prover. While their proofs are constructive, they do not support extraction of executable OCaml code. Butler et al.~\cite{butler2021formalising} formalise $\Sigma$-protocols in CryptoHOL~\cite{10.1007/978-3-662-49498-1_20}, built on Isabelle/HOL. As in~\cite{5552642}, the focus is on machine-checked correctness rather than code extraction. Moreover, the use of classical reasoning, e.g., the law of excluded middle in Isabelle/HOL may preclude extraction of executable code. Haines et al.~\cite{287095} formalise $\Sigma$-protocols in Rocq and extract OCaml code to verify the 2020 IACR election. This work is the closest to ours; however, it focuses solely on the verification of IACR elections, whereas our approach targets the correct-by-construction development of primitives for electronic voting and other privacy-preserving applications based on $\Sigma$-protocols. Their framework also lacks a natural probabilistic reasoning approach to zero-knowledge. Firsov et al.~\cite{10221929} formalise $\Sigma$-protocols in EasyCrypt (the successor to CertiCrypt), focusing on machine-checked proofs rather than executable extraction. 

  
\subsection{Limitation and Future Work}
Our formalisation lacks mix-net~\cite{10.1145/358549.358563}, a key primitive in electronic voting that breaks the link between voters and their ballots by taking encrypted ballots and outputting them in a randomised order. Since outputs remain encrypted, a malicious mix-net could replace genuine ballots with forged ones, so it must provide a zero-knowledge proof of correct shuffling~\cite{10.1007/978-3-642-02620-1_28,10.1007/978-3-642-29011-4_17}. Mix-net is also used to verify ballot validity in the Schulze voting method \cite{haines2019verifiable}. Another possibility is to formalise Bulletproofs \cite{bunz2018bulletproofs}; this would allow us to provide compact range proofs for ballot validity (e.g., proving that a vote is $0$ or $1$), and could also be adapted to construct verifiable shuffles with logarithmic proof size (we have encoded the inner-product protocol but have not yet proved any security properties about our encoding). Another promising direction is a domain-specific language (DSL) for $\Sigma$-protocols embedded in Rocq. The \texttt{\seqsplit{sigma-compiler}}~\cite{sigma-compiler,cryptoeprint:2026/794} macro for Rust allows a protocol designer to write a high-level, declarative specification---lists of public parameters, secret witnesses, and logical statements---and automatically synthesises the corresponding prover and verifier functions. While this eliminates many manual implementation errors, the generated code itself is not formally verified. By embedding a similar DSL within Rocq, one could obtain a correct-by-construction pipeline: the same specification that is used to generate executable OCaml or Rust code could also serve as the input to a formalisation that proves, inside Rocq, that the resulting proof satisfies completeness, special soundness, and honest-verifier zero-knowledge. Given our current infrastructure, it would be fairly straightforward to  encode these algorithms and prove their security properties. 
   

   Another limitation of our work is 
   that we have formalised the interactive version 
   of $\Sigma$-protocols and proved their correctness. However, 
   in reality the interaction is removed by 
   using a hash function. Ideally, one would like 
   to formalise the security of $\Sigma$-protocols in the random oracle 
   model, the closest we can get. This means that under the assumption 
   that the hash function behaves as a random oracle, the non-interactive 
   proof derived from the Fiat-Shamir transform \cite{10.5555/36664.36676} is secure. 
   The random oracle model provides a framework where we can 
   reason about the security properties  of 
   the transform. There is an existing 
   work about formalising random oracles in the Rocq theorem 
   prover \cite{10.1007/11617990_3} and we would like to 
   connect it to our implementation; however, it 
   will require some non-trivial work.
 

   In this paper, we have presented a certified implementation of a comprehensive library of $\Sigma$-protocols in the Rocq theorem prover, providing machine-checked proofs of completeness, special soundness, and special honest-verifier zero-knowledge for each construction. Our generic formalisation over abstract vector spaces captures a wide class of discrete-logarithm-based protocols while remaining independent of concrete cryptographic groups. Unlike prior work, our approach supports simultaneous extraction to OCaml and WebAssembly, eliminating code duplication between frontend and backend without sacrificing correctness. We have demonstrated the practical utility of our library by implementing a Helios election verifier (validating real IACR 2023/2024 elections) and a full client-server system for Approval Voting. The voting community can readily adopt our certified primitives to build end-to-end verifiable systems with strong, machine-checked security guarantees.



%%
%% The next two lines define the bibliography style to be used, and
%% the bibliography file.
\bibliographystyle{ACM-Reference-Format}
\bibliography{reference}


%%
%% If your work has an appendix, this is the place to put it.
\appendix
\section{Open Science}
This formalisation is implemented in Rocq~9.0.1, with additional dependencies on Dune, ExtLib, Menhir, Cryptokit, HACL$^{*}$, Lwt, Yojson Cohttp-lwt-unix. Detailed setup and usage instructions are provided in the \texttt{src/README.md}.

\end{document}
\endinput
%%
%% End of file `sample-sigplan.tex'.
